\documentclass[12pt]{article}
\usepackage[utf8]{inputenc}
\usepackage[T1]{fontenc}
\usepackage{hyperref}
\title{Maranunggu Language Resources for K-12 Curriculum Development}
\author{Compiled by Claude Sonnet 5.5 (claude-sonnet-5-5) for Andr\'as Kornai}
\date{}
\begin{document}
\maketitle

\section{Introduction}
Maranunggu (Merranunggu, Maranungku; ISO 639-3 zmr, Glottocode mara1386) is a Western Daly language of the Daly River region and the Anson Bay coast, Northern Territory, Australia. It is endangered, and only a few elderly speakers remain. Searches found no school programme, primer or classroom material for it. The resources below are mostly reference and archival.

\section{Primary Educational Resources}
No primers, textbooks or school materials for Maranunggu turned up in the searches. Any curriculum would have to be built from the reference materials below, in consultation with the community.

\section{Linguistic and Reference Materials}
\begin{itemize}
\item D.\ T.\ Tryon, \emph{An Introduction to Maranungku (Northern Australia)}, Pacific Linguistics B-15, 1970, the main grammatical description: \url{https://openresearch-repository.anu.edu.au/bitstream/1885/146471/1/PL-B15.pdf}
\item Grambank entry: \url{https://grambank.clld.org/languages/mara1386}
\item Endangered Languages Catalogue bibliography: \url{https://endangeredlanguages.com/lang/1905/bibliography}
\item Wikipedia on the Marranj dialect continuum, which includes Maranunggu: \url{https://en.wikipedia.org/wiki/Marranj_language}
\item Wikipedia on Marranunggu: \url{https://en.wikipedia.org/wiki/Marranunggu}
\item Wikipedia on the Western Daly languages: \url{https://en.wikipedia.org/wiki/Western_Daly_languages}
\end{itemize}

\section{Religious and Community Resources}
No scripture translations or community organisation pages specific to Maranunggu were found. The South Australian Museum archives list a language-group entry (as Marinunggo): \url{https://www.samuseum.sa.gov.au/collection/archives/language_groups/marinunggo}

\section{Digital and Audio Resources}
No audio, video, app or online dictionary resources were found in the searches.

\section{Curriculum Development Context}
Only about five elderly speakers remain, and the language is not known to be taught in schools. A curriculum would therefore be a documentation-based revitalisation effort. It would start from Tryon's grammar and from any archival recordings, and it should be led by the Maranunggu community. Little else is available.
\end{document}